
本章给出本书各处所用构造的细节。
先按节概述内容。

第~7.1 节给出研究各类伪全纯圆盘模空间的分析细节，
更具体地说，赋予这些空间仓西结构。
我们的目的要求在链层面工作，
并在这些模空间上构造虚基本链\liftIndex{I105}\liftIndex{I313}。
同时考虑不同 $k$ 和 $\beta$ 对应的
$\mathcal M^{\mathrm{main}}_{k+1}(\beta)$
也至关重要，
因为所得代数结构只能作为整体得到良定义。
这些模空间彼此关联：
一个模空间的边界分解为其他模空间的纤维积之并。
我们在各
$\mathcal M^{\mathrm{main}}_{k+1}(\beta)$
上构造与这一分解相容的仓西结构；
相容性的精确含义见第~7.1.1 小节。

第~7.1.2--7.1.4 小节给出
$\mathcal M^{\mathrm{main}}_{k+1}(\beta)$
仓西结构的分析构造。
第~7.1.2 小节为模空间内部的每一点构造仓西邻域；
核心是建立适于应用隐函数定理的弗雷德霍姆框架。
第~7.1.3 小节则为边界上的每一点构造仓西邻域，
所需的主要工具是博特--莫尔斯型情形中的
陶布斯型胶合论证。
第~7.1.4 小节综合前两小节的结果，
完成 $\mathcal M^{\mathrm{main}}_{k+1}(\beta)$
上仓西结构的构造。
第~7.1.5 小节讨论前面各节构造中涉及的其他模空间。
其中一个新的几何要点是定理~7.1.51：
它为定义 $A_\infty$ 映射时出现的空间给出几何构造。

本章的核心是第~7.2 节的横截性问题。
我们强调，自 1996 年虚基本链技术问世后，
横截性的关键主要在有限维拓扑，
而非无限维分析。
因此，本书把第~7.2 节的横截性问题
同第~7.1 节的弗雷德霍姆理论及胶合定理分开处理。
当前情形有两个新的困难，
需要发展超出 \textbf{[FuOn99II]} 所用方法的新技术。

第一个困难来自问题的博特--莫尔斯型性质：
不同模空间通过纤维积相联系。
若是通过直积相联系的莫尔斯型情形，
例如 \textbf{[FuOn99II]} 的定理 3.11、引理 3.14
所述虚基本链方法，
就直接给出每个有限阶段所需的横截性。
第~7.2.1--7.2.2 小节将进一步解释此点。
第二个困难是必须同时处理无穷多个模空间，
见第~7.2.3 小节。
说明这两点后，
第~7.2.4--7.2.14 小节给出横截性问题的解决细节。
为克服这两个困难，我们须大量使用 $A_\infty$ 代数语言。
由于这些小节采用的同调代数和各项证明方案
并不依赖本问题的特殊性，
同样的方法也可用于其他类似情形的横截性问题。

第~7.3 节给出同伦幺元构造的细节。

第~7.4 节给出第~3.8 节引入的算子
$\mathfrak p,\mathfrak q,\mathfrak r$ 的构造细节。
这里有两个主要问题。
一是阐明格罗莫夫--威腾不变量与算子
$\mathfrak p$ 的关系。
这同没有边界标记点的伪全纯圆盘模空间的
稳定映射“紧化”实际上并不紧有关；
第~7.4.1 小节研究这一现象。
另一问题是描述算子
$\mathfrak q,\mathfrak p,\mathfrak r$
所需的同调代数。
第~3.8 节给出刻画这些算子主要性质的公式；
第~7.4 节说明这些公式可用 $L_\infty$ 结构来解释。
这种解释不仅澄清公式的含义，
还使第~7.2 节的论证可用于构造这些算子。
为此，我们需要用同调代数语言系统地解释它们；
所需内容见第~7.4.2--7.4.11 小节。

第~7.5 节的主要目标是证明定理 V：
将系数约化到 $R$ 后，
本书的滤过 $A_\infty$ 代数与德拉姆复形同伦等价。
证明这一等价的新思路是使用连续的扰动族。

\section{仓西结构的构造}
\label{sec:7-1}

\subsection{第 7.1 节结果的陈述}
\label{subsec:7-1-1}

本节构造前面各章所用模空间的仓西结构。
仓西结构的定义见附录 A1。
各类模空间的构造几乎相同，
故第~7.1.1--7.1.4 小节集中证明命题~3.4.2；
其他情形留待第~7.1.5 小节说明。
实际上，我们将证明命题~\ref{prop:7-1-1}
和命题~\ref{prop:7-1-2}，
它们蕴涵命题~3.4.2 的精确表述。
两命题所用概念的定义亦见附录 A1。

\begin{proposition}{7.1.1}\label{prop:7-1-1}
$\mathcal M^{\mathrm{main}}_{k+1}(\beta)$
上有一个使之成为紧豪斯多夫空间的拓扑。
此外，它具有带边与角的有向仓西结构\liftIndex{I227}，
维数为 $\mu_L(\beta)+k-2+n$；
相对于此结构，评价映射
\[
\operatorname{ev}:
\mathcal M^{\mathrm{main}}_{k+1}(\beta)
\longrightarrow L^{k+1}
\]
光滑、强连续且弱浸没。
\end{proposition}

命题~\ref{prop:7-1-1} 中的仓西结构
还须满足一定的相容条件。
下面的命题~\ref{prop:7-1-2}
精确描述这些条件。
为陈述它，需要附录 A1 所讨论的
仓西结构纤维积概念。
我们比较 $S_m\mathcal M^{\mathrm{main}}_{k+1}(\beta)$
上诱导的仓西结构，
与若干 $\mathcal M^{\mathrm{main}}_{k'+1}(\beta')$
的纤维积所给的仓西结构；
在此 $\omega(\beta')<\omega(\beta)$，
或 $\omega(\beta')=\omega(\beta)$ 且 $k'<k$。

记 $S_m\mathcal M^{\mathrm{main}}_{k+1}(\beta)$
为 $\mathcal M^{\mathrm{main}}_{k+1}(\beta)$
的余维 $m$ 角层。
按式~(A1.1) 的记号，
它是
$\psi_p(S_m(V_p)\cap s_p^{-1}(0))$
的并，其中 $S_m(V_p)$ 是 $V_p$ 的余维 $m$ 角层。

取
$((\Sigma,\vec z),w)
\in S_m\mathcal M^{\mathrm{main}}_{k+1}(\beta)$。
$\Sigma$ 的不可约分支可以是圆盘或球面；
圆盘分支之并仍连通、亏格为零，
边界有 $k+1$ 个标记点。
先将此曲线分解为不可约圆盘分支，
再把球面分支构成的每棵树
附着到唯一一个不可约圆盘分支上。
由此得到
$\Sigma=\Sigma_0\cup\Sigma_1\cup\cdots\cup\Sigma_m$，
其中每个 $\Sigma_i$ 都与 $\partial\Sigma$ 相交。
构造表明，余维 $m$ 角层的这类 $\Sigma$
恰有 $m+1$ 个圆盘分支。
令
$\beta(i)=w_*[\Sigma_i]\in\Pi(M,L)$，
并设 $\Sigma_i$ 的边界有 $b_i+1$ 个奇点或标记点
（$\Pi(M,L)$ 见定义~2.4.17）。
于是 $w$ 在 $\Sigma_i$ 上的限制 $w_i$
给出
$\mathcal M^{\mathrm{main,reg}}_{b_i+1}(\beta(i))$
中的一个元素。

假设 $z_0\in\Sigma_0$，
并按定义~4.6.6 的方式排列各 $\Sigma_i$。
对每个 $i\ne0$，
存在唯一的 $j(i)$ 使
$\Sigma_{j(i)}>\Sigma_i$ 且
$\Sigma_{j(i)}\cap\Sigma_i\ne\varnothing$
（注意 $\Sigma$ 的亏格为零）。
于是交点
$\Sigma_{j(i)}\cap\Sigma_i$
是 $(\Sigma_i,w_i)$ 的一个标记点，
将 $\Sigma_i$ 的标记点重新编号，
使它成为第 $0$ 个，
记这些标记点为
$(z_0^i,\ldots,z_{b_i}^i)$。

如下将 $G_k^+$ 中的元素
$\Gamma=(T,i,v_0,V_{\mathrm{tad}},\eta)$
对应于 $((\Sigma,\vec z),w)$
（$G_k^+$ 见第~5.4.4 小节）。
$T$ 是一棵树；
内部顶点与 $\Sigma$ 的圆盘分支一一对应，
外部顶点与标记点 $z_i$ 一一对应。
两个内部顶点相连当且仅当相应圆盘分支相交；
外部顶点同内部顶点相连当且仅当
相应标记点在相应圆盘分支上。
若对应于 $\Sigma_i$ 的内部顶点 $v$
满足
$\beta(i)=w_*[\Sigma_i]=\beta_{\eta(i)}$
（$\Pi(L)$ 中第 $\eta(i)$ 个元素），
则定义 $\eta(v)=\eta(i)$。
注意 $\Sigma_i$ 包含附着于该圆盘分支的球面分支。
嵌入 $i:T\to D^2$ 由带状树\liftIndex{I103}
$T$ 的显然带状结构决定；
首个外部顶点 $v_0$ 对应第 $0$ 个标记点。
若且唯若 $\Sigma_i$ 恰有一个边界奇点、
没有标记点，对应的内部顶点属于
$V_{\mathrm{tad}}$（蝌蚪端点集）。

记 $G^+_{k,m}(\beta)$\translatornote{源书实体页 408
在这里写“$m$ 个内部顶点”。
但本页已说明余维 $m$ 角层有 $m+1$ 个圆盘分支，
且树的内部顶点与圆盘分支一一对应；
故此处计数依该对应修正，见疑误 A-0406。}
为 $G_k^+$ 中满足
$\sum_v\beta_{\eta(v)}=\beta$
且恰有 $m+1$ 个内部顶点的元素集。

定义映射
\[
\operatorname{ev}^+
=(\operatorname{ev}_1^+,\operatorname{ev}_2^+):
\prod_{i=0}^{m}
\mathcal M^{\mathrm{main}}_{b_i+1}(\beta_{\eta(i)})
\longrightarrow L^{2m}\times L^{k+1}
\]
如下。
取 $\ell(i)$ 使
$\Sigma_{j(i)}\cap\Sigma_i$
成为 $\Sigma_{j(i)}$ 的第 $\ell(i)$ 个标记点。
记
\[
\operatorname{ev}^i:
\mathcal M^{\mathrm{main}}_{b_i+1}(\beta_{\eta(i)})
\longrightarrow L^{b_i+1}
\]
为评价映射。
注意 $\Sigma$ 的边界有 $m$ 个奇点。
映射 $\operatorname{ev}^+(w_0,w_1,\ldots,w_m)$
的第一个因子定义为
\[
\begin{aligned}
\operatorname{ev}_1^+(w_0,\ldots,w_m)
=\bigl(&w_1(z_0^1),
w_{j(1)}(z_{\ell(1)}^{j(1)}),\\
&\ldots,
w_m(z_0^m),
w_{j(m)}(z_{\ell(m)}^{j(m)})\bigr).
\end{aligned}
\]

这里 $z_i^j$ 是 $\Sigma_j$ 的第 $i$ 个标记点；
$z_0^i$ 和 $z_{\ell(i)}^{j(i)}$
将成为 $\Sigma$ 上的同一点。
见图~\ref{fig:7-1-1}。

\par\medskip\noindent\begin{minipage}{\textwidth}
\centering
\setlength{\unitlength}{0.7mm}
\begin{picture}(145,76)
\linethickness{0.48pt}
\put(66,25){\circle{32}}
\put(40,40){\circle{28}}
\put(24,55){\circle{18}}
\put(94,41){\circle{32}}
\put(114,60){\circle{24}}
\put(98,68){\circle{13}}
\put(134,63){\circle{16}}
\put(30,50){\circle*{2.2}}
\put(52,33){\circle*{2.2}}
\put(80,33){\circle*{2.2}}
\put(105,52){\circle*{2.2}}
\put(103,64){\circle*{2.2}}
\put(126,62){\circle*{2.2}}
\put(66,9){\makebox(0,0){$\scriptstyle\times$}}
\put(66,41){\makebox(0,0){$\scriptstyle\times$}}
\put(49,50){\makebox(0,0){$\scriptstyle\times$}}
\put(15,55){\makebox(0,0){$\scriptstyle\times$}}
\put(114,48){\makebox(0,0){$\scriptstyle\times$}}
\put(119,50){\makebox(0,0){$\scriptstyle\times$}}
\put(62,25){$\scriptstyle\Sigma_0$}
\put(88,40){$\scriptstyle\Sigma_{j(i)}$}
\put(110,60){$\scriptstyle\Sigma_i$}
\put(62,0){$\scriptstyle z_0$}
\qbezier(105,52)(128,41)(126,21)
\put(105,12){$\scriptstyle z_0^i=z_{\ell(i)}^{j(i)}$}
\end{picture}
\par\smallskip
{\renewcommand{\thefigure}{7.1.1}
\refstepcounter{figure}
{\normalfont\bfseries 图~\thefigure}\label{fig:7-1-1}}
\end{minipage}\par\medskip

第二个因子是通常的评价映射。
具体说，若 $(\Sigma,\vec z)$ 的第 $i$ 个标记点
对应 $\Sigma_{a(i)}$ 的第 $b(i)$ 个标记点，
则
\[
\begin{aligned}
\operatorname{ev}_2^+(w_0,\ldots,w_m)
=\bigl(&w_0(z_0^0),
w_{a(1)}(z_{b(1)}^{a(1)}),\\
&\ldots,
w_{a(k)}(z_{b(k)}^{a(k)})\bigr).
\end{aligned}
\]
其中 $a(0)=b(0)=0$。
$\operatorname{ev}^+$ 由各 $\operatorname{ev}^i$ 诱导，
故由命题~\ref{prop:7-1-1} 可知，
它是光滑、强连续的映射；
该命题还蕴涵
$\operatorname{ev}_1^+$ 弱浸没。
映射 $\operatorname{ev}^+$ 由组合数据
$\Gamma=(T,i,v_0,V_{\mathrm{tad}},\eta)\in G_k^+$
决定，故记为
$\operatorname{ev}^+_\Gamma
=(\operatorname{ev}^+_{\Gamma,1},
\operatorname{ev}^+_{\Gamma,2})$。

令 $L^m\to L^{2m}$ 为对角映射
$(x_1,\ldots,x_m)
\mapsto(x_1,x_1,\ldots,x_m,x_m)$。
考虑纤维积
\[
\left(\prod_{i=0}^{m}
\mathcal M^{\mathrm{main}}_{b_i+1}
(\beta_{\eta(i)})\right)
{}_{\operatorname{ev}^+_{\Gamma,1}}
\!\times_{L^{2m}}L^m.
\]
有连续满射
\[
\bigsqcup_{\Gamma\in G^+_{k,m}(\beta)}
\left[
\left(\prod_{i=0}^{m}
\mathcal M^{\mathrm{main}}_{b_i+1}
(\beta_{\eta(i)})\right)
{}_{\operatorname{ev}^+_{\Gamma,1}}
\!\times_{L^{2m}}L^m
\right]
\longrightarrow
\overline{S}_m
\mathcal M^{\mathrm{main}}_{k+1}(\beta).
\]
目标是余维 $m$ 角层
$S_m\mathcal M^{\mathrm{main}}_{k+1}(\beta)$
的闭包。
\translatornote{源书实体页 409
在满射目标及下一句将目标模空间下标写为 $b_i+1$；
但 $b_i$ 依分支变化，
角层的底空间始终是
$\mathcal M^{\mathrm{main}}_{k+1}(\beta)$。
此处按同页命题 7.1.2 统一，见疑误 A-0407。}
此映射有限对一，
且限制在每个
\[
\left(\prod_{i=0}^{m}
\mathcal M^{\mathrm{main}}_{b_i+1}
(\beta_{\eta(i)})\right)
{}_{\operatorname{ev}^+_{\Gamma,1}}
\!\times_{L^{2m}}L^m
\quad(\Gamma\in G^+_{k,m}(\beta))
\]
上都是单射。
$\overline{S}_m\mathcal M^{\mathrm{main}}_{k+1}(\beta)$
与上述纤维积都有仓西结构。
由于评价映射被假定为弱浸没，
定义~A1.37 保证作为仓西结构的纤维积良定义。

\begin{proposition}{7.1.2}\label{prop:7-1-2}
纤维积
\[
\left(\prod_{i=0}^{m}
\mathcal M^{\mathrm{main}}_{b_i+1}
(\beta_{\eta(i)})\right)
{}_{\operatorname{ev}^+_{\Gamma,1}}
\!\times_{L^{2m}}L^m
\]
上的仓西结构，
等于
$\overline{S}_m\mathcal M^{\mathrm{main}}_{k+1}(\beta)$
上诱导的仓西结构的拉回。
\end{proposition}

本节给出的命题~\ref{prop:7-1-1}
和命题~\ref{prop:7-1-2} 的证明，
结合了 \textbf{[FOh97, FuOn99II, Fuk96II]}
中发展的技术。

